\documentclass[a4paper,12pt,twoside]{memoir}

% Castellano
\usepackage[spanish,es-tabla]{babel}
\selectlanguage{spanish}
\usepackage[utf8]{inputenc}
\usepackage[T1]{fontenc}
\usepackage{lmodern} % scalable font
\usepackage{microtype}
\usepackage{placeins}

\RequirePackage{booktabs}
\RequirePackage[table]{xcolor}
\RequirePackage{xtab}
\RequirePackage{multirow}

% Links
\PassOptionsToPackage{hyphens}{url}\usepackage[colorlinks]{hyperref}
\hypersetup{
	allcolors = {red}
}

% Ecuaciones
\usepackage{amsmath}

% Rutas de fichero / paquete
\newcommand{\ruta}[1]{{\sffamily #1}}

% Párrafos
\nonzeroparskip

% Huérfanas y viudas
\widowpenalty100000
\clubpenalty100000

% Evitar solapes en el header
\nouppercaseheads

% Imagenes
\usepackage{graphicx}
\newcommand{\imagen}[2]{
	\begin{figure}[!h]
		\centering
		\includegraphics[width=0.9\textwidth]{#1}
		\caption{#2}\label{fig:#1}
	\end{figure}
	\FloatBarrier
}

\newcommand{\imagenflotante}[3]{ %Modificado para poder introducir el tamaño de la imagen(width)
	\begin{figure}[!h]
		\centering
		\includegraphics[width=#3\textwidth]{#1}
		\caption{#2}\label{fig:#1}
	\end{figure}
	\FloatBarrier
}

% El comando \figura nos permite insertar figuras comodamente, y utilizando
% siempre el mismo formato. Los parametros son:
% 1 -> Porcentaje del ancho de página que ocupará la figura (de 0 a 1)
% 2 --> Fichero de la imagen
% 3 --> Texto a pie de imagen
% 4 --> Etiqueta (label) para referencias
% 5 --> Opciones que queramos pasarle al \includegraphics
% 6 --> Opciones de posicionamiento a pasarle a \begin{figure}
\newcommand{\figuraConPosicion}[6]{%
  \setlength{\anchoFloat}{#1\textwidth}%
  \addtolength{\anchoFloat}{-4\fboxsep}%
  \setlength{\anchoFigura}{\anchoFloat}%
  \begin{figure}[#6]
    \begin{center}%
      \Ovalbox{%
        \begin{minipage}{\anchoFloat}%
          \begin{center}%
            \includegraphics[width=\anchoFigura,#5]{#2}%
            \caption{#3}%
            \label{#4}%
          \end{center}%
        \end{minipage}
      }%
    \end{center}%
  \end{figure}%
}

%
% Comando para incluir imágenes en formato apaisado (sin marco).
\newcommand{\figuraApaisadaSinMarco}[5]{%
  \begin{figure}%
    \begin{center}%
    \includegraphics[angle=90,height=#1\textheight,#5]{#2}%
    \caption{#3}%
    \label{#4}%
    \end{center}%
  \end{figure}%
}
% Para las tablas
\newcommand{\otoprule}{\midrule [\heavyrulewidth]}
%
% Nuevo comando para tablas pequeñas (menos de una página).
\newcommand{\tablaSmall}[5]{%
 \begin{table}
  \begin{center}
   \rowcolors {2}{gray!35}{}
   \begin{tabular}{#2}
    \toprule
    #4
    \otoprule
    #5
    \bottomrule
   \end{tabular}
   \caption{#1}
   \label{tabla:#3}
  \end{center}
 \end{table}
}

%
%Para el float H de tablaSmallSinColores
\usepackage{float}

%
% Nuevo comando para tablas pequeñas (menos de una página).
\newcommand{\tablaSmallSinColores}[5]{%
 \begin{table}[H]
  \begin{center}
   \begin{tabular}{#2}
    \toprule
    #4
    \otoprule
    #5
    \bottomrule
   \end{tabular}
   \caption{#1}
   \label{tabla:#3}
  \end{center}
 \end{table}
}

\newcommand{\tablaApaisadaSmall}[5]{%
\begin{landscape}
  \begin{table}
   \begin{center}
    \rowcolors {2}{gray!35}{}
    \begin{tabular}{#2}
     \toprule
     #4
     \otoprule
     #5
     \bottomrule
    \end{tabular}
    \caption{#1}
    \label{tabla:#3}
   \end{center}
  \end{table}
\end{landscape}
}

%
% Nuevo comando para tablas grandes con cabecera y filas alternas coloreadas en gris.
\newcommand{\tabla}[6]{%
  \begin{center}
    \tablefirsthead{
      \toprule
      #5
      \otoprule
    }
    \tablehead{
      \multicolumn{#3}{l}{\small\sl continúa desde la página anterior}\\
      \toprule
      #5
      \otoprule
    }
    \tabletail{
      \hline
      \multicolumn{#3}{r}{\small\sl continúa en la página siguiente}\\
    }
    \tablelasttail{
      \hline
    }
    \bottomcaption{#1}
    \rowcolors {2}{gray!35}{}
    \begin{xtabular}{#2}
      #6
      \bottomrule
    \end{xtabular}
    \label{tabla:#4}
  \end{center}
}

%
% Nuevo comando para tablas grandes con cabecera.
\newcommand{\tablaSinColores}[6]{%
  \begin{center}
    \tablefirsthead{
      \toprule
      #5
      \otoprule
    }
    \tablehead{
      \multicolumn{#3}{l}{\small\sl continúa desde la página anterior}\\
      \toprule
      #5
      \otoprule
    }
    \tabletail{
      \hline
      \multicolumn{#3}{r}{\small\sl continúa en la página siguiente}\\
    }
    \tablelasttail{
      \hline
    }
    \bottomcaption{#1}
    \begin{xtabular}{#2}
      #6
      \bottomrule
    \end{xtabular}
    \label{tabla:#4}
  \end{center}
}

%
% Nuevo comando para tablas grandes sin cabecera.
\newcommand{\tablaSinCabecera}[5]{%
  \begin{center}
    \tablefirsthead{
      \toprule
    }
    \tablehead{
      \multicolumn{#3}{l}{\small\sl continúa desde la página anterior}\\
      \hline
    }
    \tabletail{
      \hline
      \multicolumn{#3}{r}{\small\sl continúa en la página siguiente}\\
    }
    \tablelasttail{
      \hline
    }
    \bottomcaption{#1}
  \begin{xtabular}{#2}
    #5
   \bottomrule
  \end{xtabular}
  \label{tabla:#4}
  \end{center}
}



\definecolor{cgoLight}{HTML}{EEEEEE}
\definecolor{cgoExtralight}{HTML}{FFFFFF}

%
% Nuevo comando para tablas grandes sin cabecera.
\newcommand{\tablaSinCabeceraConBandas}[5]{%
  \begin{center}
    \tablefirsthead{
      \toprule
    }
    \tablehead{
      \multicolumn{#3}{l}{\small\sl continúa desde la página anterior}\\
      \hline
    }
    \tabletail{
      \hline
      \multicolumn{#3}{r}{\small\sl continúa en la página siguiente}\\
    }
    \tablelasttail{
      \hline
    }
    \bottomcaption{#1}
    \rowcolors[]{1}{cgoExtralight}{cgoLight}

  \begin{xtabular}{#2}
    #5
   \bottomrule
  \end{xtabular}
  \label{tabla:#4}
  \end{center}
}




\graphicspath{ {./img/} }

% Capítulos
\chapterstyle{bianchi}
\newcommand{\capitulo}[2]{
	\setcounter{chapter}{#1}
	\setcounter{section}{0}
	\chapter*{#2}
	\addcontentsline{toc}{chapter}{#2}
	\markboth{#2}{#2}
}

% Apéndices
\renewcommand{\appendixname}{Apéndice}
\renewcommand*\cftappendixname{\appendixname}

\newcommand{\apendice}[1]{
	%\renewcommand{\thechapter}{A}
	\chapter{#1}
}

\renewcommand*\cftappendixname{\appendixname\ }

% Formato de portada
\makeatletter
\usepackage{xcolor}
\newcommand{\tutor}[1]{\def\@tutor{#1}}
\newcommand{\course}[1]{\def\@course{#1}}
\definecolor{cpardoBox}{HTML}{E6E6FF}
\def\maketitle{
  \null
  \thispagestyle{empty}
  % Cabecera ----------------
\noindent\includegraphics[width=\textwidth]{cabecera}\vspace{1cm}%
  \vfill
  % Título proyecto y escudo informática ----------------
  \colorbox{cpardoBox}{%
    \begin{minipage}{.8\textwidth}
      \vspace{.5cm}\Large
      \begin{center}
      \textbf{TFG del Grado en Ingeniería Informática}\vspace{.6cm}\\
      \textbf{\LARGE\@title{}}
      \end{center}
      \vspace{.2cm}
    \end{minipage}

  }%
  \hfill\begin{minipage}{.20\textwidth}
    \includegraphics[width=\textwidth]{escudoInfor}
  \end{minipage}
  \vfill
  % Datos de alumno, curso y tutores ------------------
  \begin{center}%
  {%
    \noindent\LARGE
    Presentado por \@author{}\\ 
    en Universidad de Burgos --- \@date{}\\
    Tutor: \@tutor{}\\
  }%
  \end{center}%
  \null
  \cleardoublepage
  }
\makeatother

\newcommand{\nombre}{David Renedo Gil}
\newcommand{\nombreTFG}{GII 22.24 Aplicación de gestión de TFGs}

% Datos de portada
\title{\nombreTFG \\Documentación Técnica}
\author{David Renedo Gil}
\tutor{Álvar Arnaiz González y Ana Serrano Mamolar}
\date{\today}

\begin{document}

\maketitle



\cleardoublepage



%%%%%%%%%%%%%%%%%%%%%%%%%%%%%%%%%%%%%%%%%%%%%%%%%%%%%%%%%%%%%%%%%%%%%%%%%%%%%%%%%%%%%%%%



\frontmatter


\clearpage

% Indices
\tableofcontents

\clearpage

\listoffigures

\clearpage

\listoftables

\clearpage

\mainmatter

\appendix

\apendice{Plan de Proyecto Software}

\section{Introducción}
En esta sección se detallará la planificación que se ha realizado, el estudio de viabilidad tanto de la parte económica, como temporal y de la legal. 
\section{Planificación temporal}
Se nombrarán y explicarán brevemente las tareas realizadas a lo largo del proyecto. Estas tareas se encuentran en el \href{https://github.com/drg1006/Gestor-TFG-2022}{repositorio del proyecto en Github}. 

Se añadirán gráficas para una mejor comprensión del tiempo que ha supuesto cada tarea en los (\emph{Sprints}). Esta gráfica muestra el tiempo que se tarda en comenzar y finalizar las tareas de cada uno de los ciclos.

\subsection{Sprint 0 - Puesta a punto (5/10/22 - 19/10/22)}
Puesta a punto del proyecto.
Se procederá a plantear las herramientas con las que se va a trabajar, búsqueda de alternativas y toma de contacto con las herramientas nuevas que se van a emplear.

A continuación se detallarán las tareas que se realizaron durante este primer Sprint:

\begin{itemize}
	
	\item Añadir la extensión ZenHub al navegador. 
	Desde el \textbf{Chrome Web Store} de Google Chrome se añadió la extensión \textbf{ZenHub for GitHub}. 
	
	\item Clonar en repositorio en local. Para clonarlo se ha utilizado la herramienta \textbf{Github Desktop}. Mediante en enlace \href{https://github.com/drg1006/Gestor-TFG-2022.git}{HTTP} que proporciona \emph{Github}.
	
	\item Documentación sobre Vaadin. Se procederá a estudiar el \emph{framework} Vaadin con el que se va a trabajar. A través de la página oficial de \href{https://vaadin.com/}{Vaadin} se realiza la instalación en nuestro entorno IDE \href{https://www.eclipse.org/ide/}{Eclipse} y el aprendizaje.
	
	\item Instalación JDK 11 o superior. Para utilizar la última versión de Vaadin se descargará el  \href{https://www.oracle.com/java/technologies/downloads/#jdk17-windows}{openjdk 17}. 
	
	\item Importación de un proyecto Vaadin de prueba a Eclipse. Para probar el correcto funcionamiento de Vaadin descargaremos e importaremos el proyecto de \href{https://vaadin.com/docs/latest/guide/quick-start}{prueba}.
	
	\item Clonación e imitación del repositorio en Eclipse. Trataremos de clonar e imitar el funcionamiento de la versión \href{https://github.com/dbo1001/Gestor-TFG-2021}{anterior del proyecto} sobre la que trabajamos.
	Posteriormente se descargará también el \href{https://www.oracle.com/java/technologies/downloads/#java11-windows}{openjdk 11} para tratar de clonar el repositorio que estaba en la anterior versión del proyecto. También debemos instalar la herramienta \href{https://tomcat.apache.org/}{Tomcat}.
	
	\item Comienzo de la documentación. Para ello hemos instalado las herramientas TexStudio y MikTex como se indica en \href{https://github.com/ubutfgm/plantillaLatex}{la plantilla Latex} y se ha buscado información para iniciar la documentación.
	
	\item Actualización del README.md. 
	Se modificó el README.md del proyecto para que refleje los cambios respecto a la versión anterior. 
	
	\item Búsqueda de trabajos relacionados con la gestión de TFG/TFM.
	Se realizó una investigación con el fin de encontrar proyectos similares a la aplicación web, es decir, que consistan en la gestión de trabajos de fin de grado o similares. Los proyectos encontrados serán explicados en el apartado \textbf{Trabajos relacionados} de la memoria.
	
	Se puede ver el trascurso de estas tareas en la ilustración \ref{fig:Sprint0_CONTROL}.
	
	\imagenflotante{Sprint0_CONTROL}{Gráfica Burndown Report- Sprint 1}{0.9}
\end{itemize}

\subsection{Sprint 1 - (19/10/22 - 9/11/22)}

Se procederá a estudiar el código del repositorio y a documentar el anexo.

A continuación se detallarán las tareas que se realizaron durante este primer Sprint:

\begin{itemize}
	
	\item Comienzo de la documentación del anexo. Comenzamos en este Sprint a realizar esta documentación desde TexStudio.
	
	\item Estudio del código de todos los paquetes de la carpeta src.Tanto persistence, como util, ui, security y webService.
	
	\item Se procede a buscar el error que salta al intentar ejecutar el código en local.
	
	Se puede ver el trascurso de estas tareas en la ilustración \ref{fig:Sprint1_CONTROL}.
	
	\imagenflotante{Sprint1_CONTROL}{Gráfica Burndown <<Report-- Sprint 1>>}{0.9}
\end{itemize}

\subsection{Sprint 2 - Comienzo de la programación (10/11/22 - 23/11/22)}

En este sprint se comienza a programar y añadir código principalmente arreglando bugs que existían en la versión anterior. También se investiga sobre una alternativa al uso de Heroku que ahora es de pago.

A continuación se detallarán las tareas que se realizaron durante este segundo Sprint:


\begin{itemize}
	\item Eliminación de la distinción entre mayúsculas y minúscula en los filtros. Anteriormente se tenía que introducir el nombre exacto en una columna para que se aplicase bien el filtro, ahora no existe esa distinción.
	
	\item La \emph{URL} del apartado \emph{Documentos} era errónea y se ha sustituido por la correcta.
	
	\item Actualización apartado \emph{información}. Se ha actualizado la información respecto a los tutores y la última versión.
	
	\item Investigar sobre el webscrapping. En un futuro se deberá realizar un webscrapping con la página de investigación de la ubu, por lo que se ha estudiado en qué consiste y posibles implementaciones.
	
	\item Investigación estadística errónea. La información sobre las columnas \emph{Nota, TotalDias y Repositorio} estaba mal implementada en el archivo \texttt{BaseDeDatosTFGTFM.xls} y se ha cambiado a un formato adecuado.
	
	\item Se actualiza la memoria y el anexo correspondiente al anterior Sprint.
	
	\item Elección de una alternativa a Heroku. Heroku pasa a ser de pago el 28 de noviembre de 2022, por lo que se han buscado diferentes alternativas gratuitas como la versión de \href{https://blog.heroku.com/github-student-developer-program}{Heroku para estudiantes}, una colaboración entre Heroku y \href{https://education.github.com/students}{GitHub for Students} o \href{https://northflank.com/}{Northflank}. 
	
	\item Se realizan diferentes pruebas en las plataformas para decidir cual utilizar, y finalmente se optará por usar Heroku for Students, tras desplegar el proyecto en Heroku con éxito y que al tratar de importarlo a Northflank nos indica que debemos aportar una tasa.
	
	\item Búsqueda de librerías o APIs para realizar el webscraping en nuestro proyecto. Se analizan algunas librerías cómo \href{https://jsoup.org/}{Jsoup}, \href{https://htmlunit.sourceforge.io/}{HTMLUnit} o \href{https://jaunt-api.com/}{Jaunt} y APIs como \href{https://www.octoparse.es/blog/web-scraping-api-para-extraccion-de-datos}{Octoparse}.
	
	\item Pruebas de webscraping en un entorno local para determinar cual utilizar. Se realizan una serie de pruebas(que encontramos en el apartado de Pruebas de la Documentación) y finalmente se consigue obtener el resultado que queremos mediante JSoup, por lo que será nuestra elección.
	
	Se puede ver el trascurso de estas tareas en la ilustración \ref{fig:Sprint2_CONTROL}.
	
	\imagenflotante{Sprint2_CONTROL}{Gráfica Burndown <<Report-- Sprint 2>>}{0.9}
\end{itemize}

\subsection{Sprint 3 - Implementación nuevas pantallas (23/11/22 - 14/12/22)}
En este sprint se implementará  al proyecto el proceso de \emph{webscraping} llevado a cabo en el sprint anterior, también se crearán dos nuevas pantallas.

A continuación se detallarán las tareas que se realizaron durante este segundo Sprint:

\begin{itemize}
	
	\item Guardar los datos del \emph{webscraping} en un archivo csv/xls. Se programa un código que permite guardar los resultados sacados mediante el \emph{webscraping} a los ficheros \texttt{BaseDeDatosTFGTFM.xls} y \texttt{N4 Profesores.csv}.
	
	\item Creación del \emph{mock-up} pantalla de creación de informe. Mediante \href{https://pencil.evolus.vn/}{Pencil} crearemos una vista inicial de lo que queremos que sea nuestra pantalla. Esta pantalla se utilizará para crear un informe de un determinado área a elegir por el usuario y guardará los datos en un archivo \emph{.xls}.
	
	\item Creación del \emph{mock-up} pantalla de estadísticas del profesorad. Creada también mediante \href{https://pencil.evolus.vn/}{Pencil}. Esta pantalla se utilizará para visualizar los históricos de los profesores que queramos de la \emph{EPS}, dependiendo de los departamentos y áreas que se indiquen.
	
	\item Implementación \emph{webscrap} en nuestro proyecto. Se ha introducido esta función en nuestro proyecto.
	
	\item Implementación de la pantalla Generar Informes. Se comienza y se termina de programar esta nueva pantalla. No se han creado más \emph{issues} si no que se iba comentando en esta \emph{issue} los problemas y el continuo desarrollo de la pantalla, así como las dudas planteadas.
	
	\item Implementación de la pantalla Estadísticas del Profesorado. Se comienza y se termina de programar esta nueva pantalla.
	
	\item Corrección de la memoria y anexos. Se corrige los fallos expuestos tras el \emph{feedback} del tutor Álvar Arnaiz de los ficheros latex de memoria y anexo.
	
	Se puede ver el trascurso de estas tareas en la ilustración \ref{fig:Sprint3_CONTROL}.
	
	\imagenflotante{Sprint3_CONTROL}{Gráfica Burndown <<Report-- Sprint 3>>}{0.9}
	
\end{itemize}

\subsection{Sprint 4 - Asignación TFGs (15/12/22 - 11/01/23) }

En este sprint se realizarán varias correcciones de las funcionalidades añadidas previamente y se añadirán dos nuevas pantallas, una para subir TFGs y otra para aceptar o denegarlos.

A continuación se detallarán las tareas que se realizaron durante este segundo Sprint:

\begin{itemize}
	\item Crear la pantalla para proponer TFGs siendo tutor. Se comienza y se termina de programar esta funcionalidad en este sprint tras realizar también las correcciones indicadas por los tutores en las reuniones.
	\begin{itemize}
		\item Creación del formulario.
		\item Modificación del excel.
		\item Indicar parámetros obligatorios antes de crear el TFG.
		\item Poner por defecto varios de los datos a introducir.
	\end{itemize}
	\item Corrección de los aspectos señalados por los tutores en las pestañas añadidas previamente.
	\item Documentación de memoria y anexos. Se sigue llevando a cabo la documentación de ambos informes.
	\item Creación de la pantalla de aceptación de TFGs. Esta funcionalidad permite aceptar o denegar TFGs siendo tutor.
	\begin{itemize}
		\item Crear la tabla con los proyectos con estado \emph{Pendiente}.
		\item Crear una columna con el \emph{CheckBox} para seleccionar los TFGs.
		\item Permitir cambiar el estado mediante los botones \emph{Aceptar o Denegar}.
	\end{itemize}
	\item Indicar de forma visual en que pantalla estamos, modificando el color del \emph{botón} en la barra de navegación.
	\item Actualización del despliegue de la aplicación en \emph{Heroku}, cambiando el plan del proyecto.

	Se puede ver el trascurso de estas tareas en la ilustración \ref{fig:Sprint4_CONTROL}.

	\imagenflotante{Sprint4_CONTROL}{Gráfica Burndown <<Report-- Sprint 4>>}{0.9}
\end{itemize} 

\subsection{Sprint 5 - Modificación TFGs (12/01/23 - --/--/--) }

En este sprint se realizarán varias correcciones de las funcionalidades añadidas previamente y además se modificará la pantalla de aceptación de TFGs que pasará a ser de administración, pudiendo modificar los proyectos activos.

A continuación se detallarán las tareas que se realizaron durante este segundo Sprint:

\begin{itemize}
	\item Actualización de métodos para obtener los datos de ficheros CSV. Se han implementado los métodos para obtener la información de los ficheros CSV.
	\item Implementación del nuevo sistema de \emph{Login}. Pidiendo un registro de forma opcional para asignar un rol al usuario que inicia sesión y mostrarle así las pantallas a las que puede acceder. Se ha restringido la vista a las distintas pantallas de la aplicación dependiendo de los permisos del usuario que inicia sesión.
	\item Ajustar la nueva barra de navegación ante el usuario que inicia sesión, así como del \emph{sub menu} que se crea en la pantalla de histórico tanto del profesorado como de los proyectos. Se han estudiado también otras formas para realizar la barra de navegación que finalmente no se han implementado por no ser óptimas.
	\begin{itemize}
		\item Realizar un desplegable mediante el componente \emph{MenuBar} de Vaadin, que se descartó porque no podía redimensionar.
		\item Mediante el componente \emph{Tab}. Una barra de navegación que finalmente se descartó ya que era mucho más trabajo del pensado crearla y adaptarla al proyecto.
	\end{itemize}		
	\item Creación de la cuenta de \emph{SonarCloud} para analizar la calidad del código creado y de otros proyectos añadidos, para la pantalla de \emph{Métricas}.
	\item Modificación de la pantalla de aceptación de TFGs que pasa a llamarse\emph{Administrar TFGs}. Se han añadido las siguientes funcionalidades:
	\begin{itemize}
		\item Nuevo botón de \emph{Modificación} de TFGs. Botón que solo permite seleccionar un solo TFG de la tabla.
		\item Creación de la nueva pantalla de modificación de los TFGs. Esta pantalla tiene todos los campos posibles que puede tener un proyect. Permite modificar los datos del TFG seleccionado para mantenerlo abierto y en la pantalla de activos, como para actualizarlo y cerrarlo moviéndolo a la pestaña de \emph{Históricos}.
		\item Actualización de los datos del TFG seleccionado en el archivo de la base de datos, ya sea eliminándolo de la pestaña de activos y añadiéndolo a la de históricos como actualizándolo unicamente en la pestaña de activos.
	\end{itemize}
	\item Últimas correcciones sobre la memoria y los anexos.
	\item Comentar código de los ficheros creados, tanto de los métodos cómo comentarios de ayuda durante las ejecuciones.
	\item Corrección de bugs en formatos de variables y de código.
	
\end{itemize}

\subsection{Planificación temporal}

En esta sección se mostrará una ideal inicial de como iba a ser el reparto de horas inicial y cual fue el real de las diferentes tareas durante el desarrollo del TFG. La tabla \ref{horas} muestra una predicción de la división del proyecto.

\begin{table}[]
	\label{horas}
	\centering
	\begin{tabular}{|l|c|c|}
		\hline
		\multicolumn{1}{|c|}{\textbf{TAREA}}     & \textbf{INICIAL} & \textbf{REAL} \\ \hline
		Instalación y configuración de software y hardware & 4 h & 8 h  \\ \hline
		Programación                            & 150 h  & 200 h       \\ \hline
		Documentación memoria y anexos          & 60 h   & 70 h     \\ \hline
		Estudio de las herramientas utilizadas y alternativas & 10 h & 15 h \\ \hline
		Preparación de la presentación		    & 5 h   & 8 h        \\ \hline
	\end{tabular}
	\caption{Planteamiento de horas iniciales.}
\end{table}


\section{Estudio de viabilidad}
En este apartado se detallan los costes que llevaría realizar este proyecto de forma real. Se considerarán los costes de recursos humanos, el material empleado y el \emph{Software} usado. 

\subsection{Viabilidad económica}
Primero se detallarán los costes de los recursos utilizados para la realización de este proyecto.

\begin{itemize}
	\item \textbf{Coste de personal}.
	
	En este apartado se el gasto que supone mantener a un empleado. Se calcula que se han realizado unas \emph{300} horas de trabajo durante 4 meses. 
	Con el \textit{salario medio español de un programador} 14,43 €/hora ~\cite{SalarioProgramador} y estimando unas 20 horas semanales obtenemos un salario bruto de:
	

$$ 20\frac{horas}{semana}\times14,43\frac{\text{€}}{hora}\times4\frac{semanas}{mes}=1154\text{€}\hspace{0.5em}al\hspace{0.5em}mes  $$

Este calculo hace referencia al salario bruto del empleado, para calcular el salario real del empleado debemos  obtener los impuestos que la empresa debe pagar por él. Esto se puede consultar en la página oficial de la seguridad social: \url{http://www.seg-social.es/wps/portal/wss/internet/Trabajadores/CotizacionRecaudacionTrabajadores/36537}.

Los impuestos son:
\begin{itemize}
	\tightlist
	\item 23.6\% contingencias.
	\item 5.5\% desempleo.
	\item 0.20\% FOGASA.
	\item 0.60\% formación profesional.
\end{itemize}

\imagenflotante{regimen1}{Régimen general de la Seguridad Social parte 1}{0.9}

\imagenflotante{regimen2}{Régimen general de la Seguridad Social parte 2}{0.9}

Teniendo en cuenta estos impuestos, calculamos el gasto que supone el empleado:

$$\frac{1154\frac{\text{€}}{mes}}{1-(0.236+0.055+0.002+0.006)}=1646,21\text{€}\hspace{0.5em}al\hspace{0.5em}mes$$

Tras realizar todos los cálculos,como el proyecto ha durado seis meses, el coste total será de 6584,84€.

\item \textbf{Hardware}.
El recurso hardware utilizado es un Lenovo Legion Y540 Coste aproximado: 1200€, que ya ha sido amortizado previamente, por lo que el coste real ha sido de 0 €.
 
\item \textbf{Software}.
Todas las herramientas software utilizadas han sido gratuitas a excepción del despliegue en la nube mediante \emph{Heroku}. Los créditos obtenidos sí que han sido gratuitos gracias a la conexión entre \emph{Github estudiantes y Heroku}, pero realmente supone un gasto de 7 dólares al mes.
La duración del proyecto desplegado ha sido de tres meses por lo que el coste total ha sido de 21 dólares.

 
\item \textbf{Total}.
Desestimando gastos indirectos como el internet al vivir en un piso compartido, los gastos totales son los siguientes:

\begin{table}[]
	\label{Costes totales}
	\centering
	\begin{tabular}{|l|l|}
		\hline
		\multicolumn{1}{|c|}{\textbf{Tipo de costes}}     & \textbf{Total} \\
		Personal &  6584,84\\ 
		Hardware & 0 € \\
		Software & 21 € \\	\hline
		Total & 6605,84 € \\\hline
	\end{tabular}
	\caption{Costes totales.}
\end{table}

\end{itemize}

\subsection{Viabilidad legal}
Se detallaran las licencias \emph{Software} de cada dependencia que se ha utilizado en el proyecto. En el proyecto se ha usado la misma licencia MIT que se utilizaba en la versión anterior del proyecto que permite la libre distribución del \emph{software}.

\begin{table}[]
	\label{Dependencias del proyecto}
	\centering
	\begin{tabular}{|l|l|}
		\hline
		\multicolumn{1}{|c|}{\textbf{Software}}     & \textbf{Licencia} \\ \hline
		Vaadin & Apache License 2.0 \\ \hline
		Spring Boot Maven Plugin & Apache License 2.0 \\ \hline
		JUnit & Eclipse Public License 1.0 \\	\hline
		CsvJdbc & LGPLv2 \\ \hline
		Codoid Fillo & Apache License, Version 2.0 \\ \hline
	\end{tabular}
	\caption{Licencias de las herramientas Software.}
\end{table}

También existe una cuestión de legalidad a la hora de hacer \emph{webscraping}, ya que no siempre es legal realizar este tipo de acciones sobre algunas páginas web, sobretodo si no tenemos los permisos necesarios. En nuestro caso no aplica ya que lo realizamos sobre una \emph{web} interna, por lo que no existe ningún conflicto a la hora de obtener la información.


\apendice{Especificación de Requisitos}

\section{Introducción}

Los requisitos~\cite{Requisito} o requerimientos son las características, las expectativas los aspectos esperados o las capacidades que debe cumplir el producto o servicio que generará el proyecto. Incluye además las cualidades que debe tener el proyecto para cumplir los requisitos del producto. Por esto se distingue entre requisitos funcionales y no funcionales. Los primeros definen qué debe cumplir el producto o servicio y los segundos cómo debe ser el proyecto para que el producto cumpla el propósito.
Solo se incluirán los requisitos realizados en esta mejora.

\section{Objetivos generales}
El principal objetivo de este proyecto es continuar con el desarrollo de la aplicación web respecto a la versión anterior, realizando mejoras gráficas, añadiendo nuevas pantallas y funciones.
Nos centraremos en los siguientes puntos:
\begin{itemize}
	\item Corregir \emph{bugs} o fallos encontrados en la versión previa.
	\item Incorporar la técnica \emph{webscraping} para obtener la información actualizada de los últimos tutores de la Escuela Politécnica Superior.
	\item Implantar una pantalla con un histórico de los tutores, áreas y departamentos de la universidad, atendiendo al número de TFGs asignados por curso académico. 
	\item Generar un informe para el usuario con la información de los tutores, \emph{TFGs} dirigidos, codirigidos y créditos, del área o áreas seleccionados.
	\item Implementar un sistema para que los profesores puedan realizar propuestas de \emph{TFGs} directamente desde la aplicación. Indicando todos los campos que deben tener. 
	\item Incorporar a su vez una entrada para que los administradores puedan aceptar o denegar los \emph{TFGs} que hayan sido propuestos por los tutores, además de implementar una pantalla para modificar los datos de los TFGs activos.
	\item Limitar el acceso a las pantallas dependiendo del rol del usuario que se registra en el sistema.
\end{itemize}

\section{Catálogo de requisitos}
Se describirán los requisitos específicos, funcionales y los no funcionales, exclusivamente de las nuevas funcionalidades que se han implementado en esta versión.
\subsection{Requisitos funcionales}
\begin{itemize}
		\item \texttt{RF-1  Histórico del profesorado}: mostrar el número de profesores, áreas y departamentos encontrados en la EPS.
		\begin{itemize}
			\item \texttt{RF-1.1 Parametrizar la búsqueda}: permitir al usuario escoger qué información desea obtener en la gráfica, mostrando todas las opciones de áreas, departamentos y profesores, de las que podrá seleccionar las que desee mostrar.
			\item \texttt{RF-1.2 Obtener el número de TFGs}: mostrar en la gráfica el número de TFGs por curso de los datos seleccionados.
			\item \texttt{RF-1.3 Actualizar gráfica}: darle al usuario la opción de introducir nuevos parámetros y la gráfica se actualizará \emph{clickando} en \emph{Actualizar gráfica}. A su vez, se aleatorizarán los colores que aparecen en la leyenda y las líneas de la gráfica.
		\end{itemize}
		\item \texttt{RF-2 Actualizar base de datos}: la aplicación debe obtener los datos de los tutores de la EPS.
		\begin{itemize}
				\item \texttt{RF-2.1 Preguntar al usuario}: preguntar al usuario si desea o no realizar la actualización, indicando la última fecha de modifiación y avisando de que el proceso puede tardar un tiempo.
				\item \texttt{RF-2.2 Búsqueda de datos}: los datos (nombre, apellidos y área) se obtendrán de las \emph{webs}: \href{https://investigacion.ubu.es/unidades/2682/investigadores}{Investigadores} y para obtener el departamento de cada profesor deberemos entrar en \href{https://investigacion.ubu.es/investigadores/34937/detalle}{Detalles}.
				\item \texttt{RF-2.3 Actualizar la información}: guardar los datos en el fichero \emph{BaseDeDatosTFGTFM.xls} y en el fichero \emph{N2 Profesores}, sustituyendo la información previa.				
		\end{itemize}
				
		\item \texttt{RF-3 Generar informes}: permitir al usuario seleccionar un área o áreas sobre los que generar un informe sobre los profesores de las áreas seleccionadas. En este \emph{excel} se mostrarán el número de TFGs dirigidos, codirigidos y el número de créditos asignados a cada tutor del área en el último año. El usuario deberá indicar el número de alumnos matriculados en la asignatura de \emph{TFGs}, el nombre del informe y la hoja de ruta de descarga.
		\item \texttt{RF-4 Proponer un TFG}: permitir subir un \emph{TFG} a la base de datos del sistema, introduciendo toda su información mediante \emph{el framework Vaadin}, y además se indica de forma automática que su estado está es \emph{Pendiente} de aprobación.
		\item \texttt{RF-5 Admimistrar los TFGs}: dar la capacidad de visualizar todos los \emph{TFGs} activos para poder actualizarlos.
		\begin{itemize}
			\item RF-5.1 \texttt{Aceptar/Denegar un TFG}: Permitir aprobar o denegar los proyectos que estén con estado \emph{Pendiente}. 
			\item RF-5.2 \texttt{Modificar un TFG}: Permitir actualizar al usuario la información  de un \emph{TFG}, de esta forma también se podrá cerrar un proyecto y cambiarlo a la pestaña de \emph{Histórico} o mantenerlo abierto.
		\end{itemize}
		\item \texttt{RF-6 Descargar base de datos}: se le permitirá al usuario administrador descargar la base de datos con toda la información sobre los proyectos.	
	
		
\end{itemize}

\subsection{Requisitos no funcionales}
\begin{itemize}
	\item \texttt{RNF-1 Seguridad}: la aplicación deberá verificar que la persona accediendo a las nuevas pantallas, mencionadas en los requisitos RF-4 RF-5 y RF-6, es un usuario con los privilegios requeridos.
	\item \texttt{RNF-2 Mantenimiento y escalabilidad}: permitir que la aplicación siga ordenada, no sea pesada y permita incorporar más funciones a futuro.
	\item \texttt{RNF-3 Documentación}: comentar las modificaciones realizadas para que el usuario que maneje la página sepa cómo se han realizado los procesos. 
	\item\texttt{RNF-3 Mejorar diseño}: se realizarán mejoras gráficas de la aplicación para que resulte más atractiva e informativa. Se optará por opciones intuitivas y sencillas de utilizar, pertenecientes al \emph{framework de Vaadin}.
\end{itemize}

\section{Especificación de requisitos}

\subsection{Diagrama de casos de uso}
En esta sección se mostrarán los diagramas de casos de uso. En la aplicación hay tres actores: alumno, profesor y administrador.
Este rol se le asigna una vez el usuario se registra en la aplicación en la pantalla inicial. 


\tablaSmallSinColores{Actores de la aplicación}{p{0.2\textwidth} | p{0.8\textwidth}}{Actores_aplicación}
{\textbf{Usuario} & \textbf{Funcionalidad} \\}{
	Alumno & Es el usuario que menos permisos posee, solo tiene acceso a las pantallas de información general, proyectos activos, histórico de proyectos y métricas. No tiene acceso a ninguna de las nuevas funcionalidades añadidas. \\\hline
	
	Profesor & Es el usuario con el permiso \emph{reports} para poder generar informes y proponer TFGs, también tendrá acceso a la pantalla de histórico del profesorado, además de tener acceso a todas las mencionadas en el usuario \emph{Alumno}. \\\hline
	
	Administrador & Es el usuario que más permisos tiene, con el permiso \emph{update} que le da el acceso a todas las pantallas/botones restringidos. Tiene la capacidad de administrar los \emph{TFGs}, es decir, modificar su información, aceptarlos, denegarlos, así como actualizar los datos de las bases de datos utilizadas para guardar la información. \\
}
Se puede ver un resumen de actores para los casos de uso en la imagen \ref{fig:actores}, en ella se ve en color negro las funcionalidades de la última versión y en rojo las nuevas añadidas.
\imagenflotante{actores}{Actores de la aplicación frente a las nuevas funcionalidades.}{1.3}

\tablaSmallSinColores{Caso de uso 1: Histórico del profesorado.}{p{3cm} p{.75cm} p{9cm}}{tablaCU1}{
	\multicolumn{3}{p{10.25cm}}{Caso de uso 1: Histórico del profesorado.} \\
}
{
	Descripción                            & \multicolumn{2}{p{10.25cm}}{Información estadística sobre el profesorado y una gráfica con el número de \emph{TFGs} de los parámetros seleccionados previamente.} \\\hline
	Precondiciones                         & \multicolumn{2}{p{10.25cm}}{Haber iniciado sesión y tener permisos de administrador} \\\hline
	Requisitos                         	   & \multicolumn{2}{p{10.25cm}}{Ninguna} \\\hline
	\multirow{3}{3.5cm}{Secuencia normal}  & Paso & Acción \\\cline{2-3}
	& 1    & Indicar los parámetros que el usuario quiere mostrar en la gráfica, profesores, área, departamentos.\\\cline{2-3}
	& 2    & \emph{Clickar} en el boton \emph{Actualizar gráfica}. \\\cline{2-3}
	& 3    & Repetir el proceso para cambiar los parámetros o para añadir otros tutores. \\\hline
	Postcondiciones                        & Ninguna\\\hline
	Excepciones                        & Ninguna\\\hline
	Frecuencia                             & Baja \\\hline
	Importancia                            & Media \\\hline
	Urgencia                               & Baja \\
}

\tablaSmallSinColores{Caso de uso 2:  Actualizar base de datos.}{p{3cm} p{.75cm} p{9cm}}{tablaCU2}{
	\multicolumn{3}{p{10.25cm}}{Caso de uso 2: Actualizar base de datos.} \\
}
{
	Descripción                            & \multicolumn{2}{p{10.25cm}}{Obtención de la información sobre los profesores de la EPS } \\\hline
	Precondiciones                         & \multicolumn{2}{p{10.25cm}}{Estar en la pantalla de \emph{Histórico de profesores} y comprobar que las páginas web donde se realiza la búsqueda no estén caídas} \\\hline
	Requisitos                         	   & \multicolumn{2}{p{10.25cm}}{RF-1} \\\hline
	\multirow{3}{3.5cm}{Secuencia normal}  & Paso & Acción \\\cline{2-3}
	& 1    & Preguntar al usuario si desea actualizar, el usuario pulsará el botón \emph{SI}, si quiere actualizar.\\\hline
	Postcondiciones                        & Ninguna\\\hline
	Excepciones                        & \multicolumn{2}{p{10.25cm}}{No se encuentre la ruta del archivo a actualizar o las páginas \emph{web} estén caídas.} \\\hline
	Frecuencia                             & Baja \\\hline
	Importancia                            & Alta \\\hline
	Urgencia                               & Alta \\
}

\tablaSmallSinColores{Caso de uso 3: Generar informes.}{p{3cm} p{.75cm} p{9cm}}{tablaCU3}{
	\multicolumn{3}{p{10.25cm}}{Caso de uso 3: Generar informes.} \\
}
{
	Descripción                            & \multicolumn{2}{p{10.25cm}}{ Crear un informe cuando un usuario seleccione un área o áreas sobre los que obtenerla información mencionada anteriormente de cada uno de los tutores pertenecientes a ese área. } \\\hline
	Precondiciones                         & \multicolumn{2}{p{10.25cm}}{Iniciar sesión y tener permisos de profesor o administrador} \\\hline
	Requisitos                         	   & \multicolumn{2}{p{10.25cm}}{Ninguno} \\\hline
	\multirow{3}{3.5cm}{Secuencia normal}  & Paso & Acción \\\cline{2-3}
	& 1    & Indicar el área o áreas sobre las que hacer el informe.\\\cline{2-3}
	& 2    & Indicar el número de alumnos matriculados en la asignatura TFG. \\\cline{2-3}
	& 3    & Indicar el nombre que le queremos dar al informe\\\hline{2-3}
	Postcondiciones                        & \multicolumn{2}{p{10.25cm}}{Indicar la ruta en la que descargar el informe generado} \\\hline
	Excepciones                        & \multicolumn{2}{p{10.25cm}}{Ninguna} \\\hline
	Frecuencia                             & Media \\\hline
	Importancia                            & Media \\\hline
	Urgencia                               & Alta \\
}

\tablaSmallSinColores{Caso de uso 4: Proponer un TFG.}{p{3cm} p{.75cm} p{9cm}}{tablaCU4}{
	\multicolumn{3}{p{10.25cm}}{Caso de uso 4: Proponer un TFG.} \\
}
{
	Descripción                            & \multicolumn{2}{p{10.25cm}}{ Permitir subir a la base de datos un TFG nuevo. } \\\hline
	Precondiciones                         & \multicolumn{2}{p{10.25cm}}{Ser validado como profesor o administrador a través del login} \\\hline
	Requisitos                         	   & \multicolumn{2}{p{10.25cm}}{Introducir todos los parámetros que son obligatorios.} \\\hline
	\multirow{3}{3.5cm}{Secuencia normal}  & Paso & Acción \\\cline{2-3}
	& 1    & Añadir toda la información en los campos que se indican sobre el \emph{TFG} que se quiera subir titulo, descripción, tutor1, fechaAsignación y tutor2, alumno1, alumno2 si procede.\\\cline{2-3}
	& 2    & \emph{Clickar} en el botón \emph{Crear TFG}.\\\hline
	Postcondiciones                        & \multicolumn{2}{p{10.25cm}}{Ninguna} \\\hline
	Excepciones                        & \multicolumn{2}{p{10.25cm}}{Ninguna} \\\hline
	Frecuencia                             & Media \\\hline
	Importancia                            & Alta \\\hline
	Urgencia                               & Alta \\
}

\tablaSmallSinColores{Caso de uso 5: Administrar un TFG.}{p{3cm} p{.75cm} p{9cm}}{tablaCU5}{
	\multicolumn{3}{p{10.25cm}}{Caso de uso 5: Administrar un TFG.} \\
}
{
	Descripción                            & \multicolumn{2}{p{10.25cm}}{Permitir a un administrador modificar la información de los TFGs activos además de cambiar el estado de los que están en \emph{Pendiente} } \\\hline
	Precondiciones                         & \multicolumn{2}{p{10.25cm}}{Ser validado como administrador a través del login} \\\hline
	Requisitos                         	   & \multicolumn{2}{p{10.25cm}}{Comprobar que los datos se han introducido de manera correcta y que no falte ninguno de los parámetros obligatorios.} \\\hline
	\multirow{3}{3.5cm}{Secuencia normal}  & Paso & Acción \\\cline{2-3}
	& 1    & Seleccionar el TFG/TFGs de los que se desea realizar la modificación \\\cline{2-3}
	& 2    & Realizar la operación deseada.\\\hline
	Postcondiciones                        & \multicolumn{2}{p{10.25cm}}{RF-5.1, RF-5.2} \\\hline
	Excepciones                        & \multicolumn{2}{p{10.25cm}}{Ninguna} \\\hline
	Frecuencia                             & Media \\\hline
	Importancia                            & Alta \\\hline
	Urgencia                               & Alta \\
}

\tablaSmallSinColores{Caso de uso 5.1: Aceptar/Denegar un TFG.}{p{3cm} p{.75cm} p{9cm}}{tablaCU51}{
	\multicolumn{3}{p{10.25cm}}{Caso de uso 5.1: Aceptar/Denegar un TFG.} \\
}
{
	Descripción                            & \multicolumn{2}{p{10.25cm}}{Permitir aceptar o denegar un \emph{TFG} seleccionado} \\\hline
	Precondiciones                         & \multicolumn{2}{p{10.25cm}}{Seleccionar uno o varios TFGs} \\\hline
	Requisitos                         	   & \multicolumn{2}{p{10.25cm}}{RF-5} \\\hline
	\multirow{3}{3.5cm}{Secuencia normal}  & Paso & Acción \\\cline{2-3}
	& 1    & Seleccionar el TFG/TFGs de los que se desea realizar la modificación \\\cline{2-3}
	& 2    & \emph{Clickar} en el botón de \emph{aceptar o denegar}.\\\hline
	Postcondiciones                        & \multicolumn{2}{p{10.25cm}}{Ninguna} \\\hline
	Excepciones                        & \multicolumn{2}{p{10.25cm}}{Ninguna} \\\hline
	Frecuencia                             & Media \\\hline
	Importancia                            & Alta \\\hline
	Urgencia                               & Alta \\
}
\tablaSmallSinColores{Caso de uso 5.2: Modificar un TFG.}{p{3cm} p{.75cm} p{9cm}}{tablaCU52}{
	\multicolumn{3}{p{10.25cm}}{Caso de uso 5.2: Modificar un TFG.} \\
}
{
	Descripción                            & \multicolumn{2}{p{10.25cm}}{Permitir a un administrador modificar el \emph{TFG} seleccionado en RF-5} \\\hline
	Precondiciones                         & \multicolumn{2}{p{10.25cm}}{Haber seleccionado un solo \emph{TFG} y haber \emph{clickado} en el botón \emph{Modificar}} \\\hline
	Requisitos                         	   & \multicolumn{2}{p{10.25cm}}{RF-5} \\\hline
	\multirow{3}{3.5cm}{Secuencia normal}  & Paso & Acción \\\cline{2-3}
	& 1    & Actualizar la información deseada.\\\cline{2-3}
	& 2    & \emph{Clickar} en el botón deseado: \emph{Mantener abierto o cerrarlo}\\\hline
	Postcondiciones                        & \multicolumn{2}{p{10.25cm}}{Comprobar que los datos se han introducido de manera correcta y que no falte ninguno de los parámetros obligatorios.} \\\hline
	Excepciones                        & \multicolumn{2}{p{10.25cm}}{Ninguna} \\\hline
	Frecuencia                             & Media \\\hline
	Importancia                            & Alta \\\hline
	Urgencia                               & Alta \\
}

\tablaSmallSinColores{Caso de uso 6: Descargar base de datos.}{p{3cm} p{.75cm} p{9cm}}{tablaCU6}{
\multicolumn{3}{p{10.25cm}}{Caso de uso 6:  Descargar base de datos.} \\
}
{
	Descripción                            & \multicolumn{2}{p{10.25cm}}{Permitir a un administrador descargar la base de datos utilizada en la aplicación } \\\hline
	Precondiciones                         & \multicolumn{2}{p{10.25cm}}{Ser validado como administrador a través del login} \\\hline
	Requisitos                         	   & \multicolumn{2}{p{10.25cm}}{Ninguno} \\\hline
	\multirow{3}{3.5cm}{Secuencia normal}  & Paso & Acción \\\cline{2-3}
	& 1    & \emph{Clickar en el botón} de descarga en la parte inferior de la pantalla.\\\hline
	Postcondiciones                        & \multicolumn{2}{p{10.25cm}}{Indicar la ruta de descarga} \\\hline
	Excepciones                        & \multicolumn{2}{p{10.25cm}}{Ninguna} \\\hline
	Frecuencia                             & Media \\\hline
	Importancia                            & Alta \\\hline
	Urgencia                               & Alta \\
}

\apendice{Especificación de diseño}

\section{Introducción}
En este anexo se detallarán los aspectos referentes al diseño de la aplicación en esta mejora de la aplicación.

\section{Diseño de datos}

\subsection{Ficheros de datos}

En esta nueva versión del proyecto se han añadido varios ficheros de datos nuevos con los que se va a trabajar. Para que se puedan obtener los datos de estos ficheros \textbf{se deben dar unas condiciones}.

\begin{itemize}
	\item El fichero \texttt{src/main/resources/data/BaseDeDatosTFGTFM.xls} se ha modificado, al añadir nuevos datos, con información real sobre proyectos y más peso, surgían \emph{bugs} tratando la información por lo que los campos \emph{fechaAsignacion y fechaPresentacion} en la pestaña \emph{``N3\_Historico''} se han cambiado a tipo \emph{texto} y con el formato estándar europeo DD/MM/AAAA.
	\item El campo \emph{Nota} también se ha modificado, previamente era un número aleatorio entre 5 y 10, se ha sustituido por un valor natural de tipo \emph{general}.
\end{itemize}

En el caso que no se cumpla estas condiciones en las vistas del histórico
y los proyectos activos no se mostraran de manera correcta los datos.

La nueva información obtenida mediante el \emph{webscraping} realizado se ha guardado en dos archivos.
\begin{itemize}
	\item En el archivo csv \emph{``N4\_Profesores''}, contiene en una columna separada por comas los parámetros obtenidos sobre los profesores, nombre y apellidos, área y departamentos.
	\item Esta información se ha añadido en la pestaña \emph{``N4\_Profesores''} de separados por columnas.
\end{itemize}

También se ha creado en una de las nuevas pantallas (generar informe) añadidas un fichero descargable que contiene tres columnas con la información de los profesores de los áreas seleccionadas durante el proceso de creación del proceso, todos los campos que se guardan son tipo \emph{general}.

\subsection{Diagrama de clases}

De forma general la estructura de paquetes y fichero es la misma que en el proyecto anterior \textbf{\textit{Gestor-TFG-2021}}~\cite{Gestor-TFG-2021}. Por lo que en este apartado solo se enseñará la estructura de los nuevos ficheros añadidos para las funcionalidades creadas.

\begin{itemize}
	\item En la carpeta \texttt{ubu.digit.ui.view} se han añadido cinco nuevas clases pertenecientes a las funcionalidades añadidas. Ver imagen \ref{fig:views}
	
	\imagenflotante{views}{Diagrama de clases - Vistas}{0.9}
	
	\item En la carpeta \texttt{ubu.digit.ui.entity} se ha añadido otro archivo para representar el formulario de propuesta de TFG. Ver imagen \ref{fig:entity}
	
	\imagenflotante{entity}{Diagrama de clases - Entidades}{0.9}
	
\end{itemize}

\section{Diseño procedimental}
En esta versión se ha añadido un sistema de identificación del usuario al comienzo de la aplicación, este registro le dará por lo tanto un rol al usuario. Una vez tenga asignado un rol podrá acceder a las diferentes pantallas y funcionalidades asociadas a ese actor.
El sistema sigue la siguiente lógica. Ver imagen \ref{fig:diagramaFlujo}

\imagenflotante{diagramaFlujo}{Diagrama de Flujo - Login}{0.9}

\section{Diseño arquitectónico}

El diseño arquitectónico sigue la misma estructura que encontrábamos en la versión anterior del proyecto. Con una clara distinción entre los ficheros \emph{frontend} y los ficheros \emph{backend}.
A su vez seguimos el patrón de diseño \emph{Singleton}, que es un patrón creacional que nos permite asegurarnos de que una clase tenga una única instancia, a la vez que proporciona un punto de acceso global a dicha instancia. ~\cite{Singleton}.

\apendice{Documentación técnica de programación}

\section{Introducción}
En esta sección se van a detallarlos diferentes procesos de instalación de las herramientas que se han utilizado durante el proyecto. También se especificará la estructura del proyecto, instalación de dependencias, la compilación, la ejecución del proyecto y el despliegue en Heroku. 
Algunos detalles son similares a los utilizados en el anterior proyecto de Diana, \textbf{\textit{Gestor-TFG-2021}}~\cite{Gestor-TFG-2021}.

\section{Estructura de directorios}
Se enumerarán y describirán brevemente los directorios del proyecto. Se puede encontrar el código fuente en el repositorio de Github denominado \href{https://github.com/drg1006/gestor-tfg-2022}{``Gestor-TFG-2022''}. También se indicarán los nuevos archivos añadidos en esta versión.

\begin{itemize}
	\item \texttt{/:} directorio raíz donde se ubican el README, Maven. 
	\item \texttt{/.github/workflows} los archivos de \textit{workflow} o flujo de trabajo, para la Integración continua del proyecto en GitHub.
	\item \texttt{/Documentacion} material de documentación del proyecto y prueba empleadas.
	\begin{itemize}
		\tightlist
		\item \texttt{/Documentacion/LaTeX} ficheros para generar la memoria y los anexos realizados en \emph{TexStudio}.
		\item \texttt{/Documentacion/Pruebas} aplicaciones prototipo para comenzar el aprendizaje con \href{https://vaadin.com/}{Vaadin} y pruebas realizadas con diferentes librerías durante el Web Scraping.
	\end{itemize}
	\item \texttt{/frontend} código encargado del diseño gráfico de la aplicación por el lado del cliente.
	\item \texttt{/src}, estructura de directorios \emph{backend} de la aplicación. Ya explicada en la versión anterior del proyecto:
		\begin{itemize}
			\item \texttt{/src/main/java/ubu/digit/persistence} código fuente encargado de la conexión y lectura de los ficheros de datos (fachada de datos).		
			\item \texttt{/src/main/java/ubu/digit/security} código fuente de conexión y consulta con el moodle de UbuVirtual.
			\item \texttt{/src/main/java/ubu/digit/ui} código en relación a las ventanas y vistas de la aplicación.
			\begin{itemize}
				\item \texttt{/src/main/java/ubu/digit/ui/entity} código con las entidades de los proyectos, cursos y usuarios de la aplicación.
					\begin{itemize}		
						\item \texttt{FormularioTFG.java}, nueva entidad con los campos del formulario del TFG.
					\end{itemize}
				\item \texttt{/src/main/java/ubu/digit/ui/components} código con la interfaz gráfica de la barra de navegación y la de pie de página.
				\item \texttt{/src/main/java/ubu/digit/ui/views} código con las vistas de la aplicación.
					\begin{itemize}
						\item \texttt{ProfesoresView.java}, nueva vista con la información histórica de los profesores.
						\item \texttt{ReportView.java}, nueva vista con el código para la realización de reportes.
						\item \texttt{newProjectView.java}, nueva vista con el código que permite subir propuestas de TFGs al servidor.
						\item \texttt{ManageView.java}, nueva vista con el código que permite aceptar o denegar las propuesta de TFGs con estado \emph{Pendiente}, además de permitir modificar los que están activos.
						\item \texttt{ModifyView.java}, nueva vista con el código que permite modificar los datos de un TFG seleccionado previamente en la vista de ManageView.
					\end{itemize}
				\item \texttt{/src/main/java/ubu/digit/util} incluye los métodos empleados de utilidad empleados en toda la app. 
				\item \texttt{/src/main/java/ubu/digit/webService} servicios web empleados para la consulta en moodle.
			\end{itemize}
			
			\item \texttt{/src/test} tests unitarios sobre las clases fachada “SistInfDataCsv” y “SistInfDataXls”.
			
			\item \texttt{/src/main/resources} carpeta con los datos que se van a cargar en el servidor para obtener toda la información necesaria.
		\end{itemize}
	
\end{itemize}
			
\section{Manual del programador}

A continuación se detallará el proceso de instalación de los programas
necesarios para el desarrollo de la aplicación.

\subsection{Instalación de Java}

Actualmente se sigue ejecutando con la versión de Java 11.

Para ello se debe descargar la \href{https://www.oracle.com/es/java/technologies/javase/jdk11-archive-downloads.html}{página de descargas de Oracle Java SE 11.0} y descargar la versión de JDK 11, correspondiente con el sistema operativo que se posea y su arquitectura, ya sea de 64 o 32 bits. Ver imagen \ref{fig:Descarga_JDK_11}.

Tras escoger la versión según el SO, se leerán y aceptarán las licencias de uso de Oracle \ref{fig:Descarga_JDK11_Licencia}, y se hará \emph{click} en descargar.

\imagenflotante{Descarga_JDK_11}{Descarga de JDK 11}{0.9}

También se deberá cambiar la variable de entorno de Java del sistema.

\imagenflotante{Descarga_JDK11_Licencia}{Descarga JDK 11 Licencia}{0.9}


\subsection{Instalación de Eclipse}
A continuación se instalará un entorno de desarrollo integrado(IDE) para Java, en este caso se ha utilizado \textbf{Eclipse IDE for Enterprise Java Developers} en la versión 2021-12. 

Para descargar el IDE se accederá a la \href{https://www.eclipse.org/downloads/packages/release/2021-12/r}{página de descargas de Eclipse} y descargar la opción correspondiente a nuestro sistema operativo del \textbf{Eclipse Installer 2021-12 R}. 

En el caso de los sistemas operativos Windows se descargará un archivo ejecutable que se deberá ejecutar como administrador. Una vez ejecutado se deberá seleccionar la opción ``\textbf{\textit{Eclipse IDE for Enterprise Java Developers}}''. Ver imagen \ref{fig:Eclipse_Installer}. 

\imagenflotante{Eclipse_Installer}{Seleccionar Eclipse}{0.9}

Por último seleccionaremos el JDK (11) que vayamos a utilizar y la carpeta donde queremos instalar nuestro IDE.

\subsection{Instalación del \textit{plugin de Vaadin} para Eclipse}
Una vez se haya instalado Eclipse, se procederá a añadir el plugin de Vaadin para Eclipse. Esto se realizará mediante el \textbf{Eclipse Marketplace de Eclipse} \ref{fig:Eclipse_marketplace}, el cual se encuentra en la opción de ``\textbf{\textit{Help/Eclipse Marketplace...}}'' de la barra de herramientas.

\imagenflotante{Eclipse_marketplace}{Eclipse marketplace}{0.7}

Una vez en el Eclipse Marketplace, se buscará ``\textbf{Vaadin}'' y se pulsará ``\textbf{Go}''. Tras salir el plugin ``\textbf{\textit{Vaadin Plugin for Eclipse}}'', se pulsará a ``\textbf{Install}'' y comenzará la instalación del plugin \ref{fig:Plugin_Vaadin}. En la imagen ya se muestra una vez instalado.

\imagenflotante{Plugin_Vaadin}{Plugin Vaadin}{0.9}

\section{Compilación, instalación y ejecución del proyecto}

Se explicará como compilar, instalar y ejecutar el proyecto. En el caso de
la ejecución, se detallará como hacerlo desde un terminal y mediante Eclipse (IDE).

\subsection{Descarga del repositorio}
El código fuente se encuentra en el \href{https://github.com/drg1006/Gestor-TFG-2022.git}{repositorio del proyecto} en GitHub. Para descargarlo se deberá hacer click en ``\textbf{\textit{Code}}'' y copiar la URL que aparece en el apartado de ``\textbf{HTTP}''. Con esta URL deberemos ir al ``\textbf{GitHub Desktop}'' y clonar el repositorio \ref{fig:GitHub_Code}.

\imagenflotante{GitHub_Code}{Copiar URL repositorio}{0.9}

Si se desea tener código en local se deberá descargar el zip ``\textbf{\textit{Download ZIP}}'' en la opción ``\textbf{\textit{Code}}'' anteriormente mencionada. Una vez descargado el fichero se descomprimirá y abrirá con Eclipse. 

\subsection{Compilación del proyecto}

Para compilar el proyecto en local desde terminal se usará:
\begin{itemize}
	\item Limpiar las dependencias: \texttt{mvn clean}.
	\item Instalar dependencias y compilar: \texttt{mvn install}.
	\item Instalar en modo producción (para desplegar): \texttt{mvn package -Pproduction}.  
	\item Ejecutar test: \texttt{mvn test}.	
\end{itemize}

\subsection{Ejecución del proyecto desde local}

Para la ejecución del proyecto en local desde terminal se usará:
\begin{itemize}
	\item Entrar en la terminal que utilicemos.
	\item Acceder a la carpeta donde tenemos nuestro servidor tomcat instalado y entrar en la carpeta \textbf{/bin}.
	\item Ejecutar nuestro servidor local mediante \textbf{startup} \ref{fig:tomcat_cmd}.  
	\item Entrar en el nuestro navegador en la direccion \textbf{localhost:8080}.
	\item Pulsar en la opción Manage App \ref{fig:Manage_App}.
	\item Iniciamos sesión como manager-gui.(Indicado en el archivo /conf/tomcat-users.xml).
	\item Llegaremos a la pantalla \ref{fig:Desplegar_war} y seleccionaremos el archivo .war que hemos creado al compilar nuestro proyecto con \textbf{``mvn package -Pproduction''}.
	
	\imagenflotante{tomcat_cmd}{Consola con Tomcat ejecutado}{0.9}
	
	\imagenflotante{Manage_App}{Gestor de Aplicaciones de Tomcat}{0.9}
	
	\imagenflotante{Desplegar_war}{Desplegar el archivo .war}{0.9}
	
\end{itemize}

\subsection{Ejecución del proyecto desde Eclipse IDE}

Para la ejecución del proyecto en local desde Eclipse primero debemos importar como proyecto Maven, con el \textbf{pom.xml}. Utilizaremos también un servidor local de de \href{https://tomcat.apache.org/i}{Apache Tomcat}, en concreto, la versión 9. Se puede descargar en \href{https://tomcat.apache.org/download-90.cgi}{la página oficial de Apache Tomcat}.


Una vez descargado y descomprimido, se creará un servicio de Tomcat, ver imagen \ref{fig:añadirTomcat}, con la ruta donde se tiene descargado Tomcat y se le dará un nombre, ver imagen \ref{fig:nombreServerTomcat}. Por último, se añadirá el proyecto principal ``sistinf'', ver imagen \ref{fig:tomcatAñadirProyecto}.

\imagenflotante{añadirTomcat}{Añadir servidor de Tomcat a Eclipse}{0.9}
\imagenflotante{nombreServerTomcat}{Seleccionar carpeta contenedora de Tomcat}{0.9}
\imagenflotante{tomcatAñadirProyecto}{Añadir proyectos a servidor}{0.9}

Para ejecutarlo desde Eclipse debemos también seguir todos los pasos de compilación anteriormente mencionados.

Una vez tengamos compilado nuestro código debemos ejecutarlo (click derecho en el proyecto $\rightarrow$ \texttt{Run As} $\rightarrow$ \texttt{Run on Server}).

Si no aparece la vista de los servicios se puede añadir desde la barra de herramientas$\rightarrow$ \texttt{Window} $\rightarrow$ \texttt{Show View}. Para configurar la ruta donde se ejecuta la aplicación, por defecto en \href{http://localhost:8080/}{localhost:8080/} o en ciertos casos \href{http://localhost:8080/sistinf}{localhost:8080/sistinf}.

\subsection{Alternativa a Heroku}

A partir del día 28 de noviembre Heroku dejará de ser gratuito, pero ofrecen una alternativa para estudiantes. Esta opción es un acuerdo entre el \href{https://education.github.com}{programa de estudiantes de GitHub} y Heroku. 
Esta colaboración viene explicada en la plataforma de \href{https://blog.heroku.com/github-student-developer-program}{Heroku}.

Por ello mismo migraremos nuestro proyecto a esa versión de Heroku ya que GitHub estudiantes nos proporciona una serie de créditos mensuales durante un año para desplegar nuestra aplicación.

Los pasos a seguir son los siguientes:
\begin{itemize}
	\item Obtener la verificación GitHub para estudiantes en
	 \href{https://education.github.com/students}{https://education.github.com/students}, ver imagen\ref{fig:GitHub_stu}.
	\item Seguir las instrucciones tras entrar en \emph{Sign up for Global Campus} como se indica en \ref{fig:verify_GH}.
	\item Indicar la escuela/universidad a la que pertenecemos y el uso que le vamos a dar a la cuenta, ver imagen \ref{fig:parametros_GH}.
\end{itemize}

\imagenflotante{GitHub_stu}{Página de acceso a GitHub estudiantes}{0.9}

\imagenflotante{verify_GH}{Proceso de petición de GitHub for students}{0.9}

\imagenflotante{parametros_GH}{Aportar la información necesaria para la verificación}{0.9}

Tras ser verificados por GitHub debemos conectar nuestra cuenta con Heroku desde \href{https://www.heroku.com/github-students}{https://www.heroku.com/github-students}.

Los pasos a seguir son los siguientes:
\begin{itemize}
	\item Conectar nuestra cuenta de Heroku con la cuenta de GitHub para estudiantes.
	\item Debemos añadir una tarjeta bancaria, ya que vamos a utilizar un servicio de pago pero de forma gratuita, por lo que se solicitan esos datos (que se podrán retirar en un futuro sin ningún tipo de pago realizado).
	\item Esperar a que se confirme la solicitud realizada.
	\item Podemos comprobar si hemos sido verificados si tenemos los créditos añadidos en nuestra pestaña de \emph{Account Settings} en \emph{Billing} \ref{fig:Billing}.
\end{itemize}
El resultado final debería ser el siguiente \ref{fig:final_HGH}
\imagenflotante{final_HGH}{Pestaña final tras aplicar a la oferta}{0.9}

\imagenflotante{Billing}{Créditos de nuestra cuenta}{0.9}

Una vez hemos obtenido los créditos modificamos el plan de nuestra aplicación y cambiamos al plan \emph{Basic} con un precio de 7 dólares al mes. Ver imagen \ref{fig:Basic}

\imagenflotante{Basic}{Modificar el plan de despliegue.}{0.9}

\section{Pruebas del sistema}

En esta sección se detallarán todas las pruebas realizadas en la aplicación para comprobar su correcto funcionamiento.

Se realizarán pruebas manuales para todas las nuevas funcionalidades añadidas en esta versión.

\subsection{Archivo de prueba: Prueba01}
Este archivo situado en la carpeta \texttt{Documentacion/Pruebas} denominado \texttt{Prueba01.xls}, es un archivo estructurado de igual forma que la base de datos que utilizamos de forma habitual pero con menos carga de datos.
Se cargará con la opción \emph{Actualizar base de datos} de la parte inferior de la pantalla, de acceso exclusivo para los administradores, y se probará su correcto tratamiento durante todas las pantallas de la aplicación.

Tiene un total de nueve proyectos en la pestaña de activos, tres de ellos con un estado \emph{Pendiente} y un total de 10 en la de históricos.

\begin{enumerate}
	\item \textbf{Proyectos activos}: en esta pantalla se muestran las estadísticas sobre los proyectos de la pestaña \texttt{N2 Proyecto}. Los resultados que debemos esperar son los siguientes:
		\begin{enumerate}
			\item 6 proyectos activos, ya que habiendo un total de 9, 3 de ellos son con estado \emph{Pendiente}, por lo que no se contarán.
			\item 4 proyectos sin asignar, ya que 2 de los 6 tienen alumno asignado.
			\item 17 alumnos sin proyecto asignado, obtenidos en la pestaña \texttt{N2 Alumno}
			\item 8 tutores asignados, ya que el tutor \emph{Raúl Marticorena} se repite, y solo contamos los proyectos aceptados.
			
			\item Como podemos observar en la imagen \ref{fig:P01Activos} la información estadística es correcta y la tabla muestra bien todos los proyectos activos.
		\end{enumerate}
	
	\item \textbf{Histórico de proyectos}: en esta pantalla debemos observar la información sobre los proyectos históricos. Obtendremos los siguientes datos:
			\begin{enumerate}
			\item Total de 10 proyectos y 10 alumnos asignados.
			\item Calificación media de 8.2
			\item Tiempo medio de realización: 275,6 días.
			\item 6 proyectos son del curso 2013/2014 y 4 de ellos del curso 2014/2015. 
			\item El curso 2013/2014 tiene 9 tutores asignados y el curso 2014/2015 4 profesores en total.
			
			\item En la imagen \ref{fig:P01HistProyectos} la información estadística es correcta y la tabla muestra bien todos los proyectos activos.
		\end{enumerate}
	
	\item \textbf{Histórico del profesorado}: en esta pantalla realizaremos varias pruebas para comprobar que recoge los datos correctamente. Los resultados que debemos obtener son los siguientes:
		\begin{enumerate}
			\item Para el tutor César Ignacio García tenemos dos proyectos ambos durante el curso 2013/2014 y ambos con un tutor 2 perteneciente a la EPS, por lo que cada uno suma 0.5, el resultado que debemos obtener es 1.0 en el curso 2013/2014. Observar imagen \ref{fig:P01GraficaCesar}.
			\item Para el área \emph{Lenguajes y Sistemas informáticos}. Este área tiene tutores que aparecen en la pestaña de históricos como, para el curso 2013/2014 César Ignacio, Raúl Marticorena, Carlos López y José Francisco Díez Pastor, un total de 3 proyectos dirigidos y 2 codirigidos, lo que hace una suma de 4 proyectos en total. Para el curso 2014/2015 tenemos dos tutores que pertenece a este área, Carlos López Nozal y Carlos Pardo Aguilar, que suman dos proyectos dirigidos exclusivamente por ellos. 
			
			\item Para el departamento \emph{Ingeniería Informática}, que contiene el área \emph{Ciencia de la Computación e Inteligencia Artificial} además del mencionado en el punto anterior. Obtendríamos los mismos datos y habría que sumar la participación de tutores como Álvaro Herrero Cosío, María Belén Vaquerizo García y Bruno Baruque que, en el curso 2013/2014 tienen dos participaciones como tutores de un mismo proyecto, sumando 1.0 y en el curso 2014/2015 suman un total de 1.0 por la tutoría de Álvaro en solitario.
			
			\item Los resultados de las últimas dos pruebas deberían ser los siguientes \ref{fig:P01GraficaAD}.
		\end{enumerate}

	
	\item \textbf{Creación de informes}: esta pantalla genera informes con el número de proyectos dirigidos, codirigidos y el número de créditos asignados a un tutor de un área seleccionada. Realizaremos varias pruebas para un total de 10 alumnos:
		\begin{enumerate}
			\item Área de \emph{Lenguajes y Sistemas informáticos}. De este área tenemos varios tutores mencionados en el apartado anterior, y que cuenten con proyectos activos y con alumno asignado, es decir, eliminamos los pendientes. Los resultados deben de ser los siguientes:
				\begin{enumerate}
					\item Raúl Marticorena Sánchez: 1 proyecto codirigido, 0.99 ECTS.
					\item Carlos Pardo Aguilar: 1 proyecto codirigido, 0.99 ECTS.					
					\item José Manuel Sáiz Díez: ningún proyecto activo pero es miembro del tribunal, 0.44 ECTS.
					\item Pedro Latorre Carmona: 1 proyecto codirigido, 0.99 ECTS.						
					\item José Francisco Díez Pastor: 1 proyecto codirigido, 0.99 ECTS.										
					\item Álvar Arnaiz González: ningún proyecto activo pero es miembro del tribunal, 0.44 ECTS.					
				\end{enumerate}
			\item Área de \emph{Ciencia de la Computación e Inteligencia Artificial}. Los resultados deben de ser los siguientes:
			\begin{enumerate}
				\item Ángel Arroyo Puente, ningún proyecto activo pero es miembro del tribunal, 0.44 ECTS.							
			\end{enumerate}
			\item El informe que se genere deberá ser el mismo que se puede observar en la carpeta \texttt{Documentacion/Pruebas} con el nombre \texttt{InformeP01}.
		\end{enumerate}

	
	\item \textbf{Proponer TFG}: esta pantalla permite al usuario administrador proponer proyectos a la base de datos. Se realizarán varias pruebas.
			\begin{enumerate}
				\item Parámetros por defecto: 
					\begin{enumerate}
				 		\item El último proyecto del este archivo tiene como título \emph{GII 21.04 ....}, por lo que el título de este se debería de modificar y al estar en el curso 2022/2023 debería ser 22.01.
				 		\item El tutor debe de ser el usuario que ha iniciado sesión, en este caso figura debería figurar el mio ya que me he dado permisos de administrador, en tu caso debería figurar el tuyo.
				 		\item El alumno1 por defecto es `Aalumnos sin asignar'.
				 		\item La fecha debe de ser la del día en la que se propone el trabajo.
				 		\item Se puede observar en la imagen su funcionamiento \ref{fig:P01SubirTFG}.
					\end{enumerate}
				\item Subida correcta del proyecto: una vez se ha propuesto el trabajo, debemos comprobar que ha sido correctamente generado.
					\begin{enumerate}
						\item En la pantalla de \emph{Administrar TFGs} deberá aparecer el proyecto con estado \emph{Pendiente}.
						\item Si descargamos la base de datos deberá aparecer en la pestaña \texttt{N2 Proyecto}.						
					\end{enumerate}				
			\end{enumerate}

	
	\item \textbf{Aceptar/Denegar TFG}: en esta pantalla deberán aparecer todos los proyectos de la pestaña \texttt{N2 Proyecto} de nuestra base de datos independientemente de su estado. Estado inicial en la imagen \ref{fig:P01Inicial}. Pruebas a realizar:
		\begin{enumerate}
			\item Aceptar TFG: si aceptamos un proyecto deberá cambiar su estado en esta misma pantalla y se sumará en la pestaña de activos. Ver imágenes \ref{fig:P01SelectTFG},  \ref{fig:P01AceptTFG}, \ref{fig:P01PostActivos}
			\item Denegar TFG: si denegamos un proyecto éste será eliminado de la base de datos. Ver imágenes \ref{fig:P01SelectTFG2},
			\ref{fig:P01DenegTFG}.
		\end{enumerate}

	
	\item \textbf{Modificar TFG}: para comprobar que está funcionalidad funciona debemos primero seleccionar un proyecto en la pestaña \emph{Administrar TFGs}. Realizaremos varias pruebas:
		\begin{enumerate}
			\item Modificar datos y mantener activo: primero deberemos \emph{clickar} en el botón \emph{No} en pregunta de la parte inferior de la pantalla, ver imagen \ref{fig:P01No}. A continuación, modificaremos los campos que queramos y pulsaremos en el botón \emph{Aceptar cambios y dejar abierto}. Esta información se podrá ver actualizada en la pestaña de activos si era un proyecto aceptado o en la pestaña de administración de proyectos, también cambiará en la base de datos.
			\item Modificar datos y cerrarlo: primero deberemos \emph{clickar} en el botón \emph{No} en pregunta de la parte inferior de la pantalla, ver imagen \ref{fig:P01No}, esto nos permitirá introducir información como la nota, fecha de presentación o el enlace \emph{URL} del proyecto. A continuación, modificaremos los campos que queramos y pulsaremos en el botón \emph{Aceptar cambios y mover a histórico}. Se verán los cambios de igual forma que en el apartado anterior.			
		\end{enumerate}
\end{enumerate}

\imagenflotante{P01Activos}{Comprobación pantalla activos.}{0.9}
\imagenflotante{P01HistProyectos}{Comprobación pantalla histórico de proyectos.}{0.7}
\imagenflotante{P01GraficaCesar}{Comprobación gráfica de proyectos dirigidos por César Ignacio García.}{0.7}
\imagenflotante{P01GraficaAD}{Comprobación proyectos históricos areas y departamentos.}{0.9}
\imagenflotante{P01SubirTFG}{Comprobación campos de propuesta de proyecto.}{0.9}
\imagenflotante{P01Inicial}{Estado inicial de los proyectos.}{0.9}
\imagenflotante{P01SelectTFG}{Seleccionamos un proyecto para aceptar.}{0.9}
\imagenflotante{P01AceptTFG}{Comprobación cambio de estado.}{0.9}
\imagenflotante{P01PostActivos}{Comprobamos que se ha añadido en activos.}{0.9}
\imagenflotante{P01SelectTFG2}{Seleccionamos un proyecto para denegar.}{0.9}
\imagenflotante{P01DenegTFG}{Comprobamos que se ha eliminado de la base de datos.}{0.9}
\imagenflotante{P01No}{Clickar en la opción 'No'.}{0.9}

\subsection{Archivo de prueba: Prueba02}
Al igual que en el apartado anterior realizaremos pruebas con otro excel que se encuentra en la misma ruta y contiene diferente información. Ahora solo comprobaremos los cambios que se obtienen en las estadísticas generadas por la aplicación.

Tiene un total de 48 proyectos en la pestaña de activos, cinco de ellos con un estado \emph{Pendiente} y un total de 52 en la de históricos.

\begin{enumerate}
	\item \textbf{Proyectos activos}: en esta pantalla se muestran las estadísticas sobre los proyectos de la pestaña \texttt{N2 Proyecto}. Los resultados que debemos esperar son los siguientes:
	\begin{enumerate}
		\item 43 proyectos activos, ya que habiendo un total de 48, 5 de ellos son con estado \emph{Pendiente}, por lo que no se contarán.
		\item 18 proyectos sin asignar.
		\item 6 alumnos sin proyecto asignado, obtenidos en la pestaña \texttt{N2 Alumno}.
		\item 31 tutores asignados.
		
		\item  Como podemos observar en la imagen \ref{fig:P02Activos} la información estadística es correcta y la tabla muestra bien todos los proyectos activos.
	\end{enumerate}
	
	\item \textbf{Histórico de proyectos}: obtendremos los siguientes datos:
	\begin{enumerate}
		\item Total de 52 proyectos y 52 alumnos asignados.
		\item Calificación media de 8.02
		\item Tiempo medio de realización: 284.35 días.		
		\item Tabla con estadísticas \ref{P02Historico1}.
		
		 \begin{table}[]
			\label{P02Historico1}
			\centering
			\begin{tabular}{|l|c|c|c|}
				\hline
				\multicolumn{1}{|c|}{\textbf{}} & \textbf{2014/2015} & \textbf{2015/2016} & \textbf{2018/2019}\\\hline
				Nota media & 7.67 & 8.62  & 7.89 \\ \hline
				Nº de meses & 11.40  & 6.31 & 9.18 \\ \hline
				Nº de proyectos  & 3  & 8 &19 \\ \hline
				Nº de tutores  &  5 &12  & 28  \\ \hline
			\end{tabular}
			\caption{Estadísticas del excel Prueba02. Parte 1}
		\end{table}
	
	\begin{table}[]
		\label{P02Historico2}
		\centering
		\begin{tabular}{|l|c|c|c|c|}
			\hline
			\multicolumn{1}{|c|}{\textbf{}} & \textbf{2019/2020}
			& \textbf{2020/2021}	& \textbf{2021/2022} & \textbf{2022/2023}\\\hline
			Nota media & 7.43 &7.33 & 9.50& 8.20 \\ \hline
			Nº de meses  & 11.33 & 10.18& 6.88	& 9.99 \\ \hline
			Nº de proyectos    &7  &6 & 5	&4  \\ \hline
			Nº de tutores  & 12 &11 & 7	& 8   \\ \hline
		\end{tabular}
		\caption{Estadísticas del excel Prueba02. Parte 2}
	\end{table}
		
	\end{enumerate}
	
	\item \textbf{Histórico del profesorado}:  los resultados que debemos obtener son los siguientes:
	\begin{enumerate}
		\item Tutor \emph{Raúl Marticorena Sánchez}, 2 tfgs codirigidos en el curso 2015/2016, 2 proyectos dirigidos y tres codirigidos durante el curso 2018/2019, pero en uno de ellos el tutor2 no es de la EPS, por lo que cuenta como 1 entero, y uno codirigido, también con un tutor que no es de la EPS, en el curso 2019/2020. Deberíamos obtener entonces 1 en el curso 2015/2016, 4 en el curso 2019/2020 y 1 en el curso 2020/2021. El resultado se puede ver en la imagen \ref{fig:P02GraficaRaul}.
		
		\item Tras realizar los cálculos a mano para el área de  \emph{Lenguajes y Sistemas informáticos}, se obtienen un total de 1 proyecto en 2014/2015, 4 en 2015/2016, 8 en 2018/2019, 4 en 2019/2020, 3 en 2020/2021 y 1 en 2022/2023. Los resultados deberían ser los siguientes \ref{fig:P02GraficaAD}.
	\end{enumerate}	
	
	\item \textbf{Creación de informes}: realizaremos varias pruebas para un total de 30 alumnos:
	\begin{enumerate}
		\item El área de \emph{Lenguajes y Sistemas informáticos} tiene un total de 30 proyectos asignados durante este curso.
		\item El área de \emph{Ciencia de la Computación e Inteligencia Artificial} tiene un total de 6 proyectos asignados durante este curso.
		\item El informe que se genere deberá ser el mismo que se puede observar en la carpeta \texttt{Documentacion/Pruebas} con el nombre \texttt{InformeP02}.
	\end{enumerate}
\end{enumerate}
\imagenflotante{P02Activos}{Comprobación pantalla proyectos activos.}{0.9}

\imagenflotante{P02GraficaRaul}{Comprobación gráfica de proyectos dirigidos por Raúl Marticorena Sánchez.}{0.9}
\imagenflotante{P02GraficaAD}{Comprobación gráfica de los proyectos del área de Lenguajes y Sistemas Informáticos.}{0.9}
\section{Validaciones del sistema}

En esta sección se detallarán cómo la aplicación realiza validaciones antes de
ejecutar las operaciones que el usuario desea en las nuevas pantallas añadidas.

\subsection{Creación de informes}

En esta pantalla se hacen varias comprobaciones antes de la creación del informe como se puede ver en la imagen \ref{fig:validacionInforme}.
En ella se pueden ver los campos obligatorios para la creación del informe, así como una notificación y avisos al usuario.

\imagenflotante{validacionInforme}{Validación campos del informe.}{0.9}

\subsection{Subida de TFGs}

En esta funcionalidad, al igual que en la pantalla de creación de informes, el usuario debe introducir una serie de campos obligatorios si desea proponer un proyecto. Se le avisa al usuario de los campos obligatorios como se puede ver en la imagen \ref{fig:validacionSubirTfg}.

\imagenflotante{validacionSubirTfg}{Validación campos de propuesta de TFG.}{0.9}

\subsection{Administrar TFGs}
A la hora de administrar los \emph{TFGs}, tanto para aceptar y denegar, como para modificarlos se realizan diferentes validaciones.

A la hora de y denegar se debe seleccionar al menos un proyecto como se ve en las imágenes \ref{fig:aceptTFG} y \ref{fig:denegTFG}.

\imagenflotante{aceptTFG}{Validación seleccionar TFG para aceptar.}{0.9}
\imagenflotante{denegTFG}{Validación seleccionar TFG para denegar.}{0.9}

Una vez se ha seleccionado más de un proyecto se vuelve a realizar una confirmación al usuario antes de realizar la operación como se puede ver en las imágenes \ref{fig:validAcept} y \ref{fig:validDeneg}.

\imagenflotante{validAcept}{Validación para aceptar TFGs}{0.9}
\imagenflotante{validDeneg}{Validación para denegar TFGs.}{0.9}

A la hora de seleccionar un proyecto para modificar nos salta un aviso si seleccionamos varios o si no seleccionamos ninguno como se ve en la imagen \ref{fig:validModif}

\imagenflotante{validModif}{Validación para modificar TFGs.}{0.9}

\subsection{Manejo de sesiones}

En esta versión se ha añadido la posibilidad de tener varias sesiones iniciadas a la vez en la aplicación, ya que actualmente se han añadido varios roles para los usuarios, es necesario controlar las sesiones que se generan, de esta forma no se machacan los datos de sesión de otros usuarios. Mediante el uso de variables internas se permite a diferentes usuarios iniciar sesión al mismo tiempo en la \emph{web} como se puede ver en la imagen \ref{fig:variasSesiones}, en la que en un navegador se ha iniciado sesión y en otro no.

\imagenflotante{variasSesiones}{Aplicación web con varias sesiones a la vez.}{0.9}

\subsection{Actualización de ficheros}

Aunque esta funcionalidad ya estuviese incluida en la versión anterior se ha probado que el archivo que se sube a la base de datos y el que se descarga, sin realizar ningún tipo de modificación durante su uso, es el mismo, y tiene los mismos datos e información.
\apendice{Documentación de usuario}

\section{Introducción}
A continuación se describirán los requisitos mínimos a cumplir para que el usuario pueda entrar en la aplicación y usarla.

\section{Requisitos de usuarios}
 Al estar la aplicación desplegada en \href{https://gestor-tfg-2022.herokuapp.com/}{https://gestor-tfg-2022.herokuapp.com/} por lo que solamente hará falta disponer de Internet.
\section{Instalación}
Para utilizar la aplicación no será necesario instalar ningún componente en nuestro ordenador, a excepción de un navegador web.
\section{Manual del usuario}
A continuación se detallará el uso de la web, exclusivamente de las nuevas pantallas implantadas.
\subsection{Histórico profesorado}
En esta nueva pantalla implantada se expone información histórica sobre el profesorado.

En la parte superior de la pantalla tenemos una opción para actualizar la base datos utilizada. Se informa de la última modificación de la base de datos actual y se le advierte al usuario que esta actualización es un proceso lento que puede tardar al rededor de un minuto. Ver imagen \ref{fig:ActualizarDatos}.

\imagenflotante{ActualizarDatos}{Preguntar al usuario si quiere actualizar la base de datos}{0.9}

Si se ha optado por actualizar, se mostrará al final un aviso al usuario con el tiempo transcurrido durante la operación. Ver imagen \ref{fig:Notifiacion}

\imagenflotante{Notificacion}{Notificación al usuario sobre la actualización}{0.9}

Posteriormente tenemos un pequeño apartado de información sobre el número de áreas, departamentos y profesores de la EPS. Ver imagen \ref{fig:infor_profes}.

\imagenflotante{infor_profes}{Información sobre la EPS}{0.9}

A continuación debemos escoger las áreas, departamentos y profesores que deseemos visualizar en la gráfica. Una vez seleccionados \emph{clickaremos} en actualizar gráfica.
Si queremos seleccionar varios profesores debemos introducir primero uno y posteriormente el siguiente, para eliminarlos de la selección pulsaremos en la \emph{x} del nuevo botón que se añade tras indicar un tutor. Ver imagen \ref{fig:parametros}.

\imagenflotante{parametros}{Selección de parámetros}{0.9}

La gráfica muestra el número de TFGs por curso asignado a ese parámetro. Ver imagen \ref{fig:graficaAnexo}

\imagenflotante{graficaAnexo}{Gráfica final tras seleccionar los parámetros}{0.9}

\subsection{Creación de informes}
En esta pantalla se le da al usuario la opción de añadir el número de alumnos matriculados en la asignatura \emph{Trabajos de Fin De Grado} y un área sobre el que se quiere hacer un informe, y el nombre que se le quiere dar. 
Ver imagen \ref{fig:pantalla_informes}.

\imagenflotante{pantalla_informes}{Gráfica final tras seleccionar los parámetros}{0.9}

Este informe contendrá el número total de TFGs dirigidos, codirigidos y el número de créditos asignados a cada uno de los profesores del área seleccionado en el último curso académico. Si se seleccionan varias áreas se crearán varias hojas en el documento \emph{excel} generado, con la información pertinente. 
Pestaña 1: ver imagen \ref{fig:informe_excel}
\imagenflotante{informe_excel}{Informe generado}{0.9}

Pestaña 2: ver imagen \ref{fig:informe_excel2}
\imagenflotante{informe_excel2}{Informe generado}{0.9}

\subsection{Oferta de TFG}

A esta nueva pantalla solo podrán acceder los usuarios de \emph{Moodle} que sean profesores, es decir se pedirá al usuario que se identifique con su \emph{email y contraseña}. Si no tienen los permisos necesarios se les negará el acceso.

En este nuevo menú se le pedirá al tutor que introduzca los parámetros necesarios para crear un nuevo TFG:

\begin{itemize}
	\item Titulo.
	\item Descripción.
	\item Tutor1, por defecto se indica el nombre del tutor que se ha \emph{loggeado}.
	\item Tutor2.
	\item Alumno1, por defecto se indica \emph{'Alumnos sin asignar'}.
	\item Alumno2.
	\item CursoAsignación, por defecto se indicará el curso actual.
\end{itemize}

Todos estos parámetros son modificables. Una vez se seleccionen, el usuario \emph{clickará} en \emph{Subir TFG} y se añadirá a la base de datos, en la pestaña de \emph{N2 Proyectos}, es decir, en los proyectos activos.
También se añadirá una nueva columna que indica si el TFG ha sido aceptado o si la petición está aún pendiente. En esta pestaña siempre se añadirán como \emph{pendientes} ya que en la siguiente será un administrador el que modifique su estado. Ver imagen \ref{fig:subirTFG}

\imagenflotante{subirTFG}{Pantalla para indicar la información del TFG}{0.9}

\subsection{Aceptación TFG}

A esta nueva pantalla solo podrán acceder los usuarios de \emph{Moodle} que sean administradores, es decir se pedirá al usuario que se identifique con su \emph{email y contraseña}. Si no tienen los permisos necesarios se les negará el acceso.

Aquí aparecerán una tabla con la lista de TFGs con estado \emph{Pendiente}, en una tabla con filtros. Ver imagen \ref{fig:tablaPendientes}

\imagenflotante{tablaPendientes}{Tabla con los TFGs con estado Pendiente}{0.9}

En la parte inferior aparecerán una serie de componentes de \emph{Vaadin} para actualizar el estado del TFG que se desee.

\begin{itemize}
	\item \emph{ComboBox} con los títulos del los \emph{TFGs} de la tabla superior para indicar uno. Ver imagen \ref{fig:titulosTFG}
	\item \emph{ComboBox} con las opciones de \emph{Estado} del TFG. Ver imagen \ref{fig:estadosTFG}
	\item Botón para confirmar la modificación del TFG.
\end{itemize}

\imagenflotante{titulosTFG}{Desplegable con los títulos de los TFGs}{0.9}

\imagenflotante{estadosTFG}{Desplegable con los estados del TFG}{0.9}

El \emph{ComboBox} con los estados y el botón para confirmar la modificación no están activados desde el principio, como se puede ver en la figura \ref{fig:desactivados}.

\imagenflotante{desactivados}{Componentes desactivados}{0.9}

Estos componentes se activarán a medida que seleccionemos los datos, es decir, una vez se indique el título se activará la opción de indicar el estado al que se quiere cambiar y una vez indiquemos el estado se activará el botón de actualizar la base de datos.






\bibliographystyle{plain}
\bibliography{bibliografiaAnexos}

\end{document}
